\documentclass{article}

\usepackage{course}
\usepackage{verbatim}
\usepackage{fullpage}

\def\handout{Design problems}
\def\coursenum{142}
\def\coursetitle{Program design and methodology}
\def\courseterm{Spring}
\def\courseyear{2025}
\parskip 6pt
\parindent 0pt

\begin{document}

Describe how you could solve each of the following problems using data
structures we've learned so far.
Specify the ADTs (List, Iterator, Stack, Queue) you
would use, how they would be used to support the required operations, and
how they should be implemented to optimize efficiency.
Although lists can do anything stacks or queues can, give the most
specific correct answer: if the functionality of a stack or queue is
enough, pick that.
Only answer ``list'' if neither a stack nor a queue suffices.

\begin{enumerate}

\item
You are hired to help design software to help with a key airline
operation: processing drink orders on a flight.
Once the pilot gives the ok, airline staff walk from the back of the plane
to the front taking everyone's order.
When they reach the front, they begin preparing those orders.
Drinks are prepared in groups of 10, that being the number of cups
that can fit onto a carrying tray.
The first group whose drinks are prepared are those in the front of
the plane (first class!).
The second group is the next 10 people further back in the plane and
so on.
You are in charge of developing a new application that will let
airline staff take orders and then display those orders 10 at a time
in the appropriate order.
%stack

\item
You are in charge of making the lines at airport security run
smoothly. Every day thousands of people pass through security. There
are 3 main types of people who go through this line: economy class
passengers, VIP passengers, and flight staff. There is one line but
some people have the ability to cut through. The system should allow
VIPs and flight staff go straight to the scanners. They are rare
occurrences and only appear 1 in every 100 people. 
%queue, plus something to keep track of the high-priority people (like
%a second queue)

\item
The program is a simulation of an vertical stack in an ultimate
frisbee game.\footnote{This is a simplified description and not quite
  how the game is played.} There's seven players and five of them are
in
a vertical line (cutters).
The other two (handlers) are positioned
horizontally in the back of the line (resembling the shape of an
upside-down T). One of the handlers has the disk (frisbee) and is
going to throw
it to one of the cutters. The first cutter in the line tries to get
it, and if they don't, they go back to the back of the line. This
keeps happening until one gets the disk. That cutter then dumps it to
the handler and goes to the back of the line, then it all starts
again. The handler could also throw the disk to the other handler but
they never join the line. The program should end after the line has
reset 7 times. What ADT would you use to implement this simulation and
how?
%queue of cutters

\item
You're cleaning your dorm room, but every time you start one job you
find another that needs your attention first! You try to put away your
laundry, but it's all dirty! Then when you try to wash your laundry
you find that your laundry basket is filled with junk you've been too
lazy to put away, and when you try to put some stuff away you find
some nasty old moldy food that you have to clean up, etc.
Note that each time you discover a new task, it must be completed
before the one that you were previously working on. 
What ADT would you use to help you keep track of all the subtasks you have to
do before your room is finally clean?
%stack

\end{enumerate}

\end{document}
